\section*{Análisis de los resultados}

\subsection*{Resultados reportados}

La sigueinte tabla reune los resultados de las metricas de RMSE, MAE Y MAPE de todos los modelos individuales y los compara con el modelo hibrido:

\begin{table}[H]
\centering
\caption{Comparación entre el modelo propuesto, modelos base y estudios previos (tabla 4 del artículo).}
\label{tab:full}
\begin{tabular}{lccc}
\toprule
Modelo / Estudio & RMSE & MAE & MAPE (\%) \\
\midrule
Híbrido propuesto  & 112{,}35 & 85{,}27  & 3{,}82 \\
XGBoost            & 125{,}87 & 95{,}14  & 4{,}35 \\
Pandoh et al. (2021) & 128{,}67 & 97{,}53  & 4{,}40 \\
ARIMA              & 138{,}46 & 102{,}28 & 4{,}92 \\
VAR                & 140{,}82 & 104{,}37 & 5{,}03 \\
Ridge              & 142{,}33 & 107{,}84 & 5{,}09 \\
Lasso              & 145{,}62 & 110{,}45 & 5{,}21 \\
Kong et al. (2019) & 152{,}13 & 115{,}64 & 5{,}85 \\
\bottomrule
\end{tabular}
\end{table}
El modelo hibrido supera a todos los demas modelos en todas las métricas: menor MAE (85{,}27), menor RMSE (112{,}35) y menor MAPE ($3{,}82\,\%$ ),La figura~\ref{fig:peak} muestra la evolución del pico de demanda a lo largo del período estudiado.

\begin{figure}[H]
\centering
\includegraphics[width=0.5\linewidth]{fig1_peak_demand.png}
\caption{Demanda máxima de electricidad por año (figura 1 del artículo original).}
\label{fig:peak}
\end{figure}
 Los resultados obtenidos son consistentes con las conclusiones del artículo de referencia, ya que el modelo híbrido volvió a ser la alternativa con mayor precisión frente a los modelos individuales. Aunque el orden de desempeño de ARIMA, VAR, Lasso y XGBoost difiere del reportado por Suriya y Agusthiyar (2026), debido a que el presente estudio utiliza un conjunto de datos distinto (Países Bajos en lugar de Tamil Nadu), se confirma que la integración de modelos estadísticos y de aprendizaje automático permite reducir el error de pronóstico. Esto evidencia que el enfoque híbrido posee una mayor capacidad para adaptarse a diferentes patrones de consumo eléctrico, constituyéndose como la estrategia más robusta para la predicción de la demanda máxima.
